L’architecture physique du système repose sur plusieurs composants matériels interconnectés via un réseau sécurisé. Les moniteurs Mindray et les caméras collectent les données en temps réel, qui sont ensuite transmises au système central via le réseau de surveillance interne. Les serveurs, hébergés localement ou dans le cloud, stockent et traitent les données, tandis que les appareils mobiles permettent au personnel médical de recevoir des notifications et d’accéder aux données des patients.

\begin{itemize}
    \item \textbf{Moniteurs Mindray} : Collectent les données des patients en temps réel (fréquence cardiaque, pression artérielle, saturation en oxygène, etc.) et les transmettent via des interfaces réseau ou des APIs propriétaires.
    \item \textbf{Caméras} : Capturent les écrans des machines et transmettent les images au service de reconnaissance de texte (OCR).
    \item \textbf{Serveurs} : Stockent et traitent les données collectées. Ils peuvent être hébergés localement ou dans le cloud (AWS, Azure).
    \item \textbf{Réseau de Communication} : Assure la transmission des données entre les moniteurs, les caméras, les serveurs et les applications mobiles/web.
    \item \textbf{Appareils Mobiles} : Permettent au personnel médical de recevoir des notifications push et d’accéder aux données des patients en temps réel.
\end{itemize}
